%\documentclass[handout]{ximera}
\documentclass[nooutcomes,noauthor, handout]{ximera}
\input{../../preamble}

\title{I'm the Biggest Kid!}

\begin{document}
\begin{abstract}\end{abstract}
\maketitle



\begin{problem}
    The students in Ms. Smith's class come in from recess having an argument.  Ms. Smith asks the children about the argument, and it soon becomes clear that the children can't decide who is the biggest in the class.  Why might the children be confused about this question?
    
%    \begin{solution}
%    The term ``biggest'' isn't very descriptive or very precise.  When you are measuring the attributes of the children, one child may be taller than another, but may also be skinnier.
%    \end{solution}
   
\end{problem}


\begin{problem}
    To settle the matter, Ms. Smith divides her class up into groups, and gives each group a different item to use:
    \begin{itemize}
        \item One group gets a ball of yarn.
        \item One group gets a stack of math textbooks.
        \item One group gets an unsharpened pencil.
        \item One group gets a playground ball.
        \item One group gets a yard stick (a 3-foot ruler).
    \end{itemize}
    
    How could the children in each group use this item to help them answer the question?  Be extremely specific.
    
%    \begin{solution}
%    Answers should be creative!  Find as many attributes to measure as possible.
%    \end{solution}

\end{problem}

\begin{problem}
Would you classify any of your answers to the previous problem as a comparison? As a measurement? What does your work tell you about how these ideas are related?
\end{problem}

\begin{problem}

What does your work with Ms. Smith's class tell you about how we use a unit when we measure? Be as specific as possible.
   
\end{problem}


\newpage

\begin{instructorNotes}

{\bf Main goal:} We introduce some of the concepts of measurement and get students reasoning about non-standard units. This activity is a success if the students come away with the idea that they need to choose an appropriate unit before they measure.

{\bf Overall picture:}

This activity is intended to introduce the subject of measurement.  The activity should bring up the need for measurement, the ambiguities that can occur when comparing different types of measurements, as well as give a gentle introduction to doing measurement.  We also encounter the idea of measuring different aspects of the same object, and how these different measurements can be somewhat independent.  Also included is some time to discuss children's misconceptions about measurement.

First problem:
\begin{itemize}
	\item         This question brings up the concept that the ``biggest'' is an ambiguous term.  Are the children debating who is the tallest?  Who weighs the most?  Whose arms are longest?  Have your class come up with as many ways as possible that a child could be considered the ``biggest'' - including some ways that seem a little unreasonable!
\end{itemize}


Second problem:
\begin{itemize}
 \item Here students are asked to begin to discuss measuring various aspects of the children in the class.  Most of the natural answers should be about 1-dimensional aspects of the children in the class, like their height or perhaps the distance around their middle.  Students should identify both the aspect they are measuring, as well as the unit they will use to measure (in that order).  Not all measurements need to have a unit (for instance, if the children will cut yarn to the height of each child, and then compare the lengths), but most will.  
        
\item         This is a good time to have a discussion about what a unit is, how we use it, and why choosing one is important!

\item Some of the students' suggestions will be a direct comparison between students, rather than a measurement. For example, we could cut the yarn into segments as tall as each child, and then line up the segments to see which is longest. This is a precursor to measurement, and thus an important idea! All measurement is a type of comparison, but we usually compare to a chosen unit so that we don't have to compare between objects.
        
\item         To use some of the unconventional items, students will have to be a little creative, and perhaps measure aspects of the children that are not one-dimensional.  For example, the children could try to measure the ``biggest'' child by how many math books they can hold up at once.  Another idea for the math books would be to try to figure out how many math books it would take to cover an outline of the child (the area of a child's silhouette).  The playground ball could be used to approximate the volume of a child.
        
 \item        During discussion, when students give a way to use one of the items, ask them to try to describe what, exactly, they are measuring.  Do they think that it's a length?  An area?  A volume?  None of the above?  This is also a good place to have students practice precision in their descriptions.  Make sure the entire class understands how the item is being used in measurement.

\end{itemize}

Third problem:
\begin{itemize}
       \item  Typically the students' work with measuring objects takes two forms: comparisons (for example, everyone cuts a string that's the length of their height, then we line up the strings and see which is longest) and measurements (for example, we see how many math textbooks a student can hold for one minute). You can ask this question directly or if you have everyone's ideas on the board you can say something like ``I think I see two different kinds of things happening here. Talk with your groups about whether you agree with me or not."
       \item The idea you want to bring out in this question is that we are going to investigate measurement, where instead of having to have everyone's result all together, we give everyone a number that relates them to a common unit and then we can compare these numbers. Comparison is a precursor to measurement, and so it's good to understand, but we will work primarily with measurement after this.
       \item This idea is very similar to one-to-one correspondence vs counting. You can bring this idea up if students don't notice it for themselves.
\end{itemize}

Fourth problem:
\begin{itemize}
        \item We would like students to start to recognize that they are making copies of a unit when they do their measurement. We'll expand this into the four-step process of measurement soon, but seeing the copies now is a good outcome for this activity.
        \item You should also discuss the dimension issues with using a unit here. If we are measuring something one-dimensional, we want to use a one-dimensional unit. This is the first time we want to encourage students to be incredibly specific when they describe their unit of measure (for example, the length of the pencil - 1D - versus the pencil - 3D.
\end{itemize}




Take-aways from this activity: the concept of ``biggest" doesn't have much meaning unless we are more specific.  Draw out, but not yet answer or address many issues in measurement.  For example, what is an appropriate unit to use when measuring?  These ideas will be addressed in future activities - it's okay to not feel ``finished" with the discussion here.

{\bf Good language:} The more specific students can be with their ideas, the better. If they can make drawings or sketches of their ideas, this can help them communicate. This activity is also a good reminder of our work on dimension. For instance, if they are using the diameter of the playground ball as a length unit, we don't want them to say only ``I used the ball" but rather ``I used the distance all the way across the ball'' or something similar. A 1D aspect must be measured with a 1D unit.

{\bf Suggested timing:} This activity and its discussion should take the entire class period.  Give students about 5 minutes to discuss the situation in Problem 1 amongst themselves, and then discuss with the class for about 5-10 minutes.  Let students think about problems 2-4 for the next 20 minutes, and use the rest of the period to discuss their ideas.

\end{instructorNotes}



\end{document}